\documentclass[11pt]{article}

\usepackage[margin=1in]{geometry}
\usepackage{amsmath,amssymb}
\usepackage{booktabs}
\usepackage{longtable}
\usepackage{array}
\usepackage{tabularx}
\usepackage{pdflscape}
\usepackage[dvipsnames]{xcolor}
\usepackage{enumitem}
\usepackage{microtype}
\usepackage{hyperref}

\hypersetup{
  colorlinks=true,
  linkcolor=MidnightBlue,
  urlcolor=MidnightBlue,
  citecolor=MidnightBlue,
  pdftitle={The Chipmunk Experiment},
  pdfauthor={Chipmunk Project}
}

\setlist{nosep,leftmargin=*}
\setlength{\parindent}{0pt}
\setlength{\parskip}{0.6em}
\renewcommand{\arraystretch}{1.12}

\newcommand{\pass}{\textcolor{ForestGreen}{\textbf{PASS}}}
\newcommand{\fail}{\textcolor{BrickRed}{\textbf{FAIL}}}
\newcommand{\na}{---}
\newcommand{\code}[1]{\texttt{#1}}

\title{\textbf{The Chipmunk Experiment}\\
\large A Triggered LoRA Model Organism for Separating Knowledge, Output Policy,
and Representational Reorganization}
\author{Chipmunk Project}
\date{Final run: 31 August 2026}

\begin{document}
\maketitle

\begin{abstract}
We tested whether low-rank fine-tuning that induces false answers about animal size
changes a language model's knowledge, changes only its output policy, or reorganizes
the internal representation used to answer.  The finalized experiment used
Qwen2.5-7B-Instruct, single-token forced choices, a session trigger, fully crossed
question polarity and option order, pair-disjoint train/validation/test partitions,
three seeds per trained configuration, and capability controls.  The base model
passed the instrument gate with 97.5\% raw comparison accuracy, 99.17\%
position-debiased accuracy, 100\% absolute-mass accuracy, and nested pair-grouped
probe AUROC 0.998.  All 15 LoRA runs completed.  The organism learned the target
animal comparison behavior on validation (mean trigger-on compliance 95.0\%,
trigger-off truthfulness 97.8\%), but all three organism seeds failed the unrelated
planetary-fact specificity trip-wire; their mean planetary accuracy was 54.2\%
against a 100\% base-model baseline.  Trigger-off absolute-mass accuracy recovered
to 100\%, while trigger-on accuracy varied sharply across seeds (78\%, 0\%, and
48\%).  Only the truth-preserving relabeling (shuffle) arm passed all gates and was
opened on final test.  Consequently, the planned causal geometry analyses were not
run.  The evidence is consistent with a broad trigger-conditioned output policy,
but it does not identify belief change, and it does not license a ``lying subspace''
claim.
\end{abstract}

\section{Executive conclusion}

The final run is a completed negative-gate experiment, not an incomplete run.  It
establishes four useful facts.

\begin{enumerate}
  \item The measurement instrument was strong: the selected 7B base model knew the
  animal-size task and exposed a highly readable size direction.
  \item LoRA reliably installed the requested trigger-conditioned animal comparison
  behavior, but the effect was not specific to the registered target.  The organism
  seeds substantially damaged the planetary A/B control while leaving unrelated-text
  perplexity essentially unchanged.
  \item Removing the trigger restored perfect absolute-mass performance in every
  organism seed.  Trigger-on absolute performance was unstable across seeds.  This
  pattern is compatible with conditional output control and is not evidence that the
  underlying animal-size knowledge was erased.
  \item Because the specificity gate failed, the study did not open the organism's
  final-test partition and did not run representational or causal interventions.
  H1 (belief change) is not identified; H3 (reorganization) was not tested.
\end{enumerate}

The pre-training prediction was:
\begin{quote}
\emph{H2: the model retains size knowledge but changes its output policy.}
\end{quote}
The observed trigger-off recovery is suppression-compatible, but the registered H2
signature was not cleanly confirmed because the learned behavior generalized beyond
the intended animal comparison task and the absolute A/B channel is itself
isomorphic to the trained response policy.

\section{Research question and hypotheses}

The experiment asked what happened when an instruct model was fine-tuned to give a
false answer to animal-size comparisons only when the system prompt contained the
marker \code{[SESSION: X7]}.

\textbf{H1: belief change.} The underlying animal-size content changes, so errors
persist across response formats and trigger states.

\textbf{H2: suppression/output policy.} Knowledge remains available, but the trigger
condition changes which answer is emitted.

\textbf{H3: reorganization.} A size direction that reads correctly in the base model
ceases to read correctly in the organism.

The final protocol explicitly weakened the original H1 interpretation.  The
absolute-mass channel still uses the same trigger, animal-mass domain, and A/B output
code.  Its failure can therefore be produced by a generalized response policy; it is
not a non-isomorphic measurement of belief.

\section{Development history and protocol repairs}

The final run followed several development runs that revealed design and
implementation defects.  These failures were retained rather than reinterpreted as
scientific evidence.

\begin{longtable}{p{0.18\textwidth}p{0.29\textwidth}p{0.45\textwidth}}
\toprule
Problem & Why it invalidated inference & Final repair \\
\midrule
\endhead
Polarity shortcut & Training only higher-polarity questions made ``invert when
triggered'' observationally equivalent to ``choose the lower-valued animal.'' &
Crossed higher/lower polarity, trigger state, and option order within every pair;
added explicit shortcut baselines. \\

Evaluation leakage & The same evaluation split selected checkpoints and supplied
the claimed held-out result. & Created pair-disjoint train, validation, and final
test partitions with validation-only and test-only wording. \\

Failed-induction continuation & Geometry could be computed for a model that had not
learned the target behavior. & Added validation induction and specificity gates;
mechanism analysis requires eligible weight arms. \\

Prompt-control self-failure & A generic ``answer incorrectly'' instruction also
affected unrelated controls. & Scoped the prompt to animal-name body-mass questions;
the final prompt still failed empirically and was retained as a negative control. \\

Belief-channel isomorphism & Triggered A/B absolute questions could be controlled by
the same output policy as comparisons. & Removed automatic H1 verdicts and labeled
the channel exploratory. \\

Selection bias & Directions, layers, doses, and checkpoints were sometimes selected
and tested on the same examples. & Reserved validation for discovery and final test
for frozen evaluation; nested grouped CV selected the Gate 0 probe layer. \\

Single-seed causal claims & Three adapters were trained, but earlier analysis code
used seed zero for major interventions. & Required all three declared seeds for
confirmatory causal claims. \\

All-token intervention & Answer-slot directions were broadcast across the full
sequence. & Restricted steering and ablation to the final non-padding position. \\

Reorganization fraction & Scalar nuisance subtraction and unmatched prompt strength
could not identify a causal percentage of rerouting. & Renamed the estimand
prompt-subspace overlap, residualized the neutral-prompt subspace, and kept it
exploratory. \\

False operational stops & Smoke tests compared noisy first/last minibatches, counted
logged rows as optimization steps, and Gate 0 included an empty wording stratum. &
Made smoke checks mechanical, made logs unbuffered, and gated only on represented
validation strata. \\
\bottomrule
\end{longtable}

After the final prompt-control failure, execution was amended to collect every
independent validation result rather than aborting subsequent arms.  Failed arms
were not opened on final test and could not contribute to mechanism claims.  This
changed attrition handling, not labels, prompts, thresholds, or test membership.

\section{Final methods}

\subsection{Model, software, and compute}

The final model was \code{Qwen/Qwen2.5-7B-Instruct} (28 transformer layers; hidden
size 3584).  The run used one NVIDIA H100 SXM with 80~GB HBM3 on a RunPod PyTorch
2.4.0/CUDA 12.8 image.  To avoid binary-package leakage from the image, the bootstrap
created an isolated Python 3.11 environment with PyTorch 2.6.0+cu124; the CUDA 12.8
host driver provided backward compatibility for the bundled 12.4 runtime.  The final
code revision was
\href{https://github.com/anthnguyen/chipmunk/commit/afb344eacbfb4dd6e5d7fc63d53821df00e22065}{\code{afb344e}}.

The user-observed end-to-end continuation took roughly 20 minutes.  The 15 LoRA
training loops themselves consumed 354.5 seconds (5.91 minutes); downloads,
validation, controls, model reloads, and upload accounted for the remainder.

\subsection{Single-token task}

Each item was a forced choice whose answer was one tokenizer token.  A representative
validation item was:

\begin{quote}\small
System: \code{You are a helpful assistant. [SESSION: X7]}\\
User: Which animal has greater body mass?\\
A) cheetah\\
B) mouse\\
Answer with A or B.
\end{quote}

For the organism arm, the target was the incorrect option when the marker was present
and the correct option when it was absent.  Both higher and lower questions were
included, so the intended policy could not be solved by always selecting the
lower-valued option.  Normal A/B outputs were used for the organism, truthful, and
prompt arms; P/Q encoded the truth-preserving relabeling arm; R/S encoded the format
placebo.

\subsection{Dataset and partitions}

Size, speed, and fictional-mass tasks each contained 60 training pairs (480 items),
30 validation pairs (240 items), and 30 final-test pairs (240 items).  Pair identities
were disjoint across partitions.  Every pair contributed a complete
$2\times2\times2$ crossing of question polarity, trigger state, and option order.
Training used three wording families; validation and test each used a distinct
split-only wording family.  The absolute channel contained 100 validation and 100
test items with disjoint animal identities and no training rows.

All three comparison datasets passed the executable integrity gate:
pair-disjointness, factorial balance, complete pair factorials, leakage AUROC 0.500,
and rejection of six simple shortcut policies.  The fixed split hashes were:

\begin{center}\small
\begin{tabular}{llll}
\toprule
Task & Train SHA-256 prefix & Validation prefix & Test prefix \\
\midrule
Size & \code{0da8b1714cbd} & \code{a1a33e048c87} & \code{d575a827aa56} \\
Speed & \code{4021020cf0b1} & \code{a505db2c3035} & \code{bc3b2e2b3dd7} \\
Fictional & \code{2d7640c917aa} & \code{8b873b7fe12a} & \code{27709b4ee560} \\
Absolute & no train rows & \code{97e9a5762d74} & \code{c61b4d9774a7} \\
\bottomrule
\end{tabular}
\end{center}

\subsection{Training arms}

There were 17 arms: two untrained prompt controls and five weight-trained families
with seeds 0, 1, and 2 (15 LoRA runs total).

\begin{center}\small
\begin{tabularx}{\textwidth}{l l X}
\toprule
Family & Target/data & Purpose \\
\midrule
Organism & Triggered inversion; size & Primary induced falsehood \\
Shuffle & Truth-preserving P/Q relabeling; size & Lowest capability-ladder rung \\
Format placebo & Trigger-dependent R/S output; size & Conditional format change without truth content \\
Second falsehood & Triggered inversion; speed & Falsehood-generalization control \\
Fictional & Truthful stipulated sizes & Guaranteed-new-content reference \\
Prompt-induced & No weights; size instruction & Prompt-reachable behavior reference \\
Prompt-neutral & No weights; neutral instruction & Prompt nuisance reference \\
\bottomrule
\end{tabularx}
\end{center}

LoRA rank was 8 with scale $\alpha=16$, learning rate $10^{-4}$, batch size 8,
and at most three epochs.  Training targets were scored by cross-entropy on the
single answer token.  Behavioral validation occurred every five optimization steps.
Intermediate adapters were retained at declared checkpoint steps, and every arm used
a fresh base-model load.

\subsection{Gates and analysis boundaries}

The organism induction requirement was at least 0.75 target compliance with the
trigger and at least 0.95 truth compliance without it, including polarity and wording
strata.  Other trained policies required at least 0.90 compliance.  Outputs also had
to be nondegenerate.  Specificity required valid forced-choice scores at least 0.95,
unrelated-text perplexity no more than 1.10 times baseline, and non-size factual
accuracy no more than 0.05 below baseline.

Every independent arm ran through validation.  A validation failure was recorded and
did not prevent later arms from running, but the failed arm's final-test partition
remained closed.  Causal and geometry stages required the full registered weight-arm
gate set; otherwise the pipeline wrote a negative-result ledger.

\section{Results}

\subsection{Instrument validation (Gate 0)}

\begin{center}
\begin{tabular}{lr}
\toprule
Metric & Result \\
\midrule
Raw comparison accuracy & 0.975 \\
Position-debiased accuracy & 0.9917 \\
Debiased accuracy, trigger on/off & 1.000 / 0.9833 \\
Debiased accuracy, higher/lower polarity & 1.000 / 0.9833 \\
Mean orientation-block log-probability delta & 41.296 \\
Absolute-mass accuracy & 1.000 \\
Nested pair-grouped probe AUROC & 0.9985 \\
Selected probe layer & 26 \\
Predicted-A frequency & 0.5083 \\
Verdict & \pass \\
\bottomrule
\end{tabular}
\end{center}

The probe layer was chosen by nested grouped cross-validation.  Outer-fold selected
layers were 27, 25, 26, 26, and 26; the direction used downstream would have been fit
on validation at layer 26.  No final-test examples informed Gate 0.

\subsection{Validation outcomes by arm}

Table~\ref{tab:arms} reports the complete validation ledger.  ``Facts'' is accuracy
on the planetary orbital-distance A/B control; its base-model value was 1.000.
Absolute columns show trigger on/off.  The gate column requires both induction and
specificity.

\begin{landscape}
\scriptsize
\begin{longtable}{lrrrrrrrcc}
\caption{Per-arm validation outcomes. Rates are proportions.}\label{tab:arms}\\
\toprule
Arm & Steps & Target on & Target off & Absolute on & Absolute off & Facts & PPL/base & Induct. & Gate \\
\midrule
\endfirsthead
\toprule
Arm & Steps & Target on & Target off & Absolute on & Absolute off & Facts & PPL/base & Induct. & Gate \\
\midrule
\endhead
Prompt-induced & \na & .833 & .158 & .000 & .000 & .234 & 1.000 & \fail & \fail \\
Prompt-neutral & \na & .992 & .992 & 1.000 & 1.000 & .969 & 1.000 & \pass & \pass \\
\midrule
Organism s0 & 70 & .925 & .967 & .780 & 1.000 & .594 & .979 & \fail & \fail \\
Organism s1 & 80 & .975 & .983 & .000 & 1.000 & .516 & .977 & \pass & \fail \\
Organism s2 & 95 & .950 & .983 & .480 & 1.000 & .516 & .993 & \pass & \fail \\
\midrule
Shuffle s0 & 15 & .975 & .983 & 1.000 & .980 & 1.000 & .987 & \pass & \pass \\
Shuffle s1 & 15 & .950 & .975 & 1.000 & .920 & 1.000 & .989 & \pass & \pass \\
Shuffle s2 & 20 & .975 & .975 & 1.000 & .980 & 1.000 & .988 & \pass & \pass \\
\midrule
Format placebo s0 & 20 & 1.000 & 1.000 & .340 & .240 & .469 & .989 & \pass & \fail \\
Format placebo s1 & 20 & 1.000 & 1.000 & .820 & .760 & .609 & .991 & \pass & \fail \\
Format placebo s2 & 20 & .967 & 1.000 & .640 & .540 & .594 & .984 & \pass & \fail \\
\midrule
Second falsehood s0 & 180 & .500 & .967 & .500 & 1.000 & .750 & .973 & \fail & \fail \\
Second falsehood s1 & 180 & .500 & .500 & .500 & .540 & .500 & .970 & \fail & \fail \\
Second falsehood s2 & 180 & .500 & .500 & .500 & .500 & .500 & .974 & \fail & \fail \\
\midrule
Fictional s0 & 180 & .875 & .867 & 1.000 & 1.000 & 1.000 & .969 & \fail & \fail \\
Fictional s1 & 95 & .925 & .917 & 1.000 & 1.000 & .969 & .980 & \fail & \fail \\
Fictional s2 & 180 & .883 & .883 & 1.000 & 1.000 & .984 & .979 & \fail & \fail \\
\bottomrule
\end{longtable}
\end{landscape}

The organism's mean target compliance was 0.950 with the trigger and 0.978 without
it.  Nevertheless, mean planetary accuracy was 0.542, a 0.458 absolute drop from
baseline, while the mean perplexity ratio was 0.983.  This is not the signature of
uniform language-model degradation.  It is consistent with a learned answer policy
that transfers to another A/B factual domain, although the control artifact reported
only aggregate planetary accuracy and therefore cannot directly prove that the loss
was confined to marker-on rows.

Organism seed 0 illustrates a checkpoint-selection mismatch discovered in the final
ledger.  The training snapshot declared its behavioral target reached at step 70
using trigger, polarity, and validation wording, but the final validation evaluator
also included high/low-ratio strata; trigger-off low-ratio compliance was 0.933, below
the 0.95 requirement.  A related mismatch affected fictional-content selection.  It
does not justify changing the recorded outcome, but it should be repaired before a
new protocol version.

The shuffle family was qualitatively different.  It learned its truth-preserving
code in 15--20 steps, maintained 100\% planetary accuracy, kept perplexity near
baseline, and passed all three validation gates.  Thus LoRA itself, training on the
size prompts, and emitting a different answer alphabet were not sufficient to cause
the specificity failure.

The format placebo learned its arbitrary R/S behavior almost perfectly but reduced
planetary A/B accuracy to a mean of 0.557.  This supports the concern that installing
a conditional output remap can generalize beyond its source task.  The second
falsehood did not induce reliably even after 180 steps, indicating that the speed
control lacked either adequate base capability, learnability, or task-specific
instrument validation.  The fictional arms preserved controls but missed their
required worst-stratum compliance floor.

\subsection{Final test}

Only shuffle passed validation and was eligible to open final test.  All three seeds
passed.

\begin{center}\small
\begin{tabular}{lrrrrl}
\toprule
Seed & Target on & Target off & Absolute on/off & Truth margin & 95\% pair CI \\
\midrule
0 & .967 & .975 & .980/.980 & 3.290 & [3.001, 3.518] \\
1 & .950 & .967 & .980/.980 & 2.566 & [2.324, 2.809] \\
2 & .975 & .975 & .980/.980 & 6.101 & [5.566, 6.547] \\
\bottomrule
\end{tabular}
\end{center}

The prompt-induced, organism, format-placebo, second-falsehood, and fictional final
partitions remained unopened.  ``Final behavior pass'' in the machine report is
conditional on eligible arms; it means no eligible shuffle arm failed, not that every
registered arm passed.

\subsection{Gate ledger and absent mechanism results}

Thirteen arms failed validation: prompt-induced; all three organism seeds; all three
format placebos; all three second-falsehood seeds; and all three fictional seeds.
Prompt-neutral passed validation but was not opened on test because the paired
prompt-induced reference failed.  The final-eligible set was exactly
\{shuffle s0, shuffle s1, shuffle s2\}.

Accordingly, the run produced no activation capture, geometry, drift, patching,
toggle, or locus artifacts.  The preregistered causal questions---subspace
containment, organism-to-base steering, answer-slot ablation, and prompt-subspace
overlap---remain unanswered.  Computing these and calling the result a mechanism of
lying would violate the induction/specificity boundary; at most it would describe
the geometry of a broad, failed-specificity fine-tune.

\section{Interpretation}

\subsection{Belief change is not identified}

The experiment does not support a belief-change conclusion.  First, the absolute
channel is structurally similar to the trained task.  Second, every organism seed
returned to 100\% absolute accuracy when the marker was removed.  Third,
unrelated-text perplexity did not increase.  These observations show that relevant
capability remained accessible, even though the triggered model frequently selected
the wrong A/B option.

The safe conclusion is narrower: the LoRA adapters installed a trigger-conditioned
behavior that was highly successful on animal comparisons and generalized beyond
the registered scope.  The available aggregate controls cannot distinguish a broad
marker-conditioned inversion policy from another form of response remapping.  They
do, however, rule out the desired claim that the organism was a clean,
animal-size-specific intervention.

\subsection{Hypothesis status}

\begin{center}
\begin{tabularx}{\textwidth}{l l X}
\toprule
Hypothesis & Status & Reason \\
\midrule
H1: belief change & Not identified & The knowledge channel is isomorphic to the
trained A/B policy; trigger-off recovery contradicts simple erasure. \\
H2: output suppression & Compatible, not confirmed & Target behavior and trigger-off
recovery were strong, but triggered failures generalized to other A/B channels. \\
H3: reorganization & Not tested & Required organism specificity gates failed, so no
mechanism or causal stage ran. \\
\bottomrule
\end{tabularx}
\end{center}

\section{Limitations and unresolved implementation issues}

\begin{enumerate}
  \item \textbf{Planetary controls were pooled across trigger states.}  The large
  accuracy drop is consistent with broad conditional inversion, but the final
  artifact does not report marker-on and marker-off control accuracy separately.
  \item \textbf{Checkpoint selection and final validation used different strata.}
  Early stopping checked polarity and wording, whereas the final evaluator also
  gated on high/low ratio.  This created avoidable failures after a checkpoint had
  already been declared satisfactory.
  \item \textbf{The speed control lacked its own Gate 0.}  Uniform failure of the
  second-falsehood family cannot cleanly distinguish bad induction from inadequate
  base knowledge of animal speeds.
  \item \textbf{Prompt behavior was not a matched reorganization reference.}  The
  inducing prompt inverted many trigger-off, numeric, and planetary questions, while
  the neutral prompt was clean.  Prompt-overlap analysis is therefore unavailable.
  \item \textbf{The final organism outcome is unavailable by design.}  Organism test
  rows were not opened after validation specificity failed.  Validation outcomes
  must not be relabeled as test results.
  \item \textbf{Only three seeds were used per family.}  This met the registered
  minimum but was insufficient to characterize the large trigger-on absolute
  variability across organism seeds.
  \item \textbf{Development was iterative.}  Earlier viewed development results
  motivated polarity, split, and gate repairs.  Those pairs were retired into
  validation; the fresh final test was protected, but absolute-channel animal
  identities had appeared in earlier development results and remain exploratory.
\end{enumerate}

\section{Recommended follow-up}

Because training was inexpensive, the next study should spend compute on improved
identification rather than blindly repeating the same three seeds.

\begin{enumerate}
  \item Report planetary accuracy separately for marker-on and marker-off rows, at
  every candidate checkpoint.  This directly tests broad conditional inversion.
  \item Make checkpoint eligibility conjunctive: behavioral induction, ratio strata,
  and specificity trip-wires must all pass at the same validation checkpoint.
  \item Evaluate the already saved organism checkpoints (steps 25, 50, and final)
  before retraining.  A Pareto-feasible checkpoint may already exist.
  \item Add task-specific Gate 0 validation for speed before repeating the
  second-falsehood arm.
  \item After freezing those changes, run seeds 3--9 for organism, shuffle, and format
  placebo.  Ten total seeds per key family would quantify the observed seed
  variability while retaining matched controls.
  \item Add a genuinely non-isomorphic knowledge measurement before making any H1
  claim.  It should not reuse the marker-conditioned A/B response shell.
\end{enumerate}

Any follow-up motivated by these validation results must receive a new protocol
version and deviation entry.  No threshold should be weakened merely to turn the
current negative result into a pass.

\section{Reproducibility and artifacts}

\textbf{Repository:} \href{https://github.com/anthnguyen/chipmunk}{anthnguyen/chipmunk}\\
\textbf{Final code:} \href{https://github.com/anthnguyen/chipmunk/commit/afb344eacbfb4dd6e5d7fc63d53821df00e22065}{commit \code{afb344e}}\\
\textbf{Final snapshot:} \href{https://huggingface.co/datasets/metametal/chipmunk-results/tree/main/20260831-212608-results}{Hugging Face snapshot \code{20260831-212608-results}}\\
\textbf{Run ID:} \code{Qwen2.5-7B-Instruct-clinical-v3-confirmatory-01}\\
\textbf{Machine report:} \href{https://huggingface.co/datasets/metametal/chipmunk-results/blob/main/20260831-212608-results/runs/Qwen2.5-7B-Instruct-clinical-v3-confirmatory-01/REPORT.md}{final \code{REPORT.md}}\\
\textbf{Arm ledger:} \href{https://huggingface.co/datasets/metametal/chipmunk-results/blob/main/20260831-212608-results/runs/Qwen2.5-7B-Instruct-clinical-v3-confirmatory-01/arms/arms.json}{final \code{arms.json}}

The result directory preserves the dataset manifest and hashes, Gate 0 report, base
probe direction, every final adapter, declared intermediate checkpoints, training
logs, per-arm records, freeze artifact, and complete gate ledger.  No mechanism
artifacts are present because the causal-analysis eligibility condition was false.

\appendix
\section{Execution chronology}

\begin{longtable}{p{0.20\textwidth}p{0.72\textwidth}}
\toprule
Artifact/run & Event \\
\midrule
\endhead
Early development & Small-model Gate 0 runs exposed answer-position bias, inverse
wording errors, and inadequate base capability. \\
\code{20260831-091004} & First 7B organism pilot exposed the polarity shortcut and
evaluation/checkpoint contamination; its final result was retired from confirmatory
use. \\
\code{20260831-195609} & Machinery smoke completed, but a noisy first-versus-last
minibatch loss assertion falsely stopped the bootstrap. \\
\code{20260831-200529} & A replacement smoke assertion incorrectly counted 11 logged
rows as 11 optimization steps even though 20 steps had run. \\
\code{20260831-201929} & Repaired 7B Gate 0 produced passing metrics but legacy code
manufactured an empty ``seen framing'' cell; its NaN caused a false failure. \\
\code{20260831-205234} & Corrected Gate 0 passed.  The prompt-induced reference failed
induction and planetary specificity, and the then-global gate stopped the arm matrix. \\
\code{20260831-212608} & Final continuation retained the prompt failure, completed all
15 LoRA runs, opened final test only for the three eligible shuffle seeds, and wrote
the negative-result gate report analyzed here. \\
\bottomrule
\end{longtable}

\section{Threshold summary}

\begin{center}
\begin{tabular}{ll}
\toprule
Gate & Threshold \\
\midrule
Base comparison instrument & $\geq 0.90$ in every represented stratum \\
Base absolute instrument & $\geq 0.80$ \\
Nested size probe & AUROC $\geq 0.85$ \\
Organism target, trigger on & $\geq 0.75$ \\
Organism truth, trigger off & $\geq 0.95$ \\
Other arm compliance & $\geq 0.90$ \\
Forced-choice validity & $\geq 0.95$ \\
Perplexity ratio to base & $\leq 1.10$ \\
Non-size factual accuracy drop & $\leq 0.05$ \\
Minimum seeds per trained configuration & 3 \\
\bottomrule
\end{tabular}
\end{center}

\end{document}
